% SASYR 2025 - Template based on:
% samplepaper.tex, a sample chapter demonstrating the
% LLNCS macro package for Springer Computer Science proceedings;
% Version 2.20 of 2017/10/04
%
\documentclass[runningheads]{llncs}
%
\usepackage{graphicx}
\usepackage{url}
\usepackage{booktabs} % For professional table formatting

\begin{document}
%
% ==========================================
% i. COVER PAGE INFO
% ==========================================
\title{Software Design Document (SDD): The Game}
\subtitle{A 2D Side-Scrolling Platformer Engine Specification}

\author{Lead Designer\inst{1} \and
Physics Programmer\inst{2} \and
Level Designer\inst{3}}

\authorrunning{L. Designer et al.}

\institute{Core Systems Department, Pixel Studio, Portugal\\
\email{designer@pixelstudio.pt}
\and
Gameplay Mechanics Lab, Tech Institute, Portugal\\
\email{physics@techinst.pt}
\and
World Building Division, Creative Arts Academy, Portugal\\
\email{levels@creativearts.pt}
}

\maketitle              % typeset the header of the contribution

\begin{abstract}
This document serves as the primary Software Design Document (SDD) for the 2D side-scrolling game project \textit{The Game}. It outlines the application architecture, kinematic equations, control mapping pipelines, database designs, and system component interfaces required for cross-platform compliance. 

\keywords{Software Design Document \and Game Architecture \and SQLite \and Kinematics.}
\end{abstract}

% ==========================================
% ii. TABLE OF CONTENTS
% ==========================================
% Note: Springer LLNCS style generates a TOC automatically if required,
% but we explicitly map the requested structure here.
\addtocontents{toc}{\protect\contentsline{section}{1 Introduction}{1}}
\addtocontents{toc}{\protect\contentsline{section}{2 System Architecture}{2}}
\addtocontents{toc}{\protect\contentsline{section}{3 User Interface}{3}}
\addtocontents{toc}{\protect\contentsline{section}{4 Glossary}{4}}
\addtocontents{toc}{\protect\contentsline{section}{5 References}{5}}

% ==========================================
% iii. VERSIONS TABLE
% ==========================================
\section*{Versions Table}
\begin{table}[h]
\centering
\caption{Document Revision History}
\label{tab:versions}
\begin{tabular}{|p{1.5cm}|p{2.5cm}|p{5.5cm}|p{2.5cm}|}
\hline
\textbf{Version} & \textbf{Date} & \textbf{Description} & \textbf{Author} \\ \hline
1.0 & May 15, 2026 & Initial Snapshot: Core physics and systems. & Project Lead \\ \hline
\end{tabular}
\end{table}

\newpage

% ==========================================
% iv. SECTIONS
% ==========================================

\section{Introduction}

\subsection{Purpose}
This document provides a comprehensive technical overview of \textit{The Game}, a 2D side-scrolling platformer. It serves as the primary technical baseline for developers, artists, and quality assurance testers to ensure consistency throughout the implementation of the game.

\subsection{Intended Audience}
The target audience includes software engineers implementing physics algorithms, level designers building game spaces via data arrays, and QA technicians validating kinematical integrity against defined thresholds.

\subsection{Overview}
\textit{The Game} focuses on high-momentum platforming mechanics. Players navigate grid-based environments using an array of physics-defying abilities. The foundational design paradigm emphasizes a decoupled architecture to let mechanics iterate independently of rendering updates.


\section{System Architecture}

\subsection{Workflow Execution}
The application executes within a deterministic, fixed-timestep game loop ticking at a strict 60Hz. The execution order for every frame is defined sequentially as follows:
\begin{enumerate}
    \item \textbf{Input Capture:} Read hardware register maps (Keyboard/Controller).
    \item \textbf{State Evaluation:} Process transitions inside the character Finite State Machine (FSM).
    \item \textbf{Physics Integration:} Apply gravity vectors and impulse velocity parameters.
    \item \textbf{Collision Resolution:} Query structural layouts and shift coordinates outwards.
    \item \textbf{Rendering Pipeline:} Generate the graphic output matrix.
\end{enumerate}

\subsection{Module Breakdown}
The application splits core behaviors into isolation barriers:
\begin{itemize}
    \item \textbf{Kinematic Controller:} Resolves motion velocity math without standard physics engine jitter.
    \item \textbf{Environmental Collision Module:} Inspects targeted tile layouts to report collision hits.
    \item \textbf{Event Manager:} Dispatches visual particle feedback and dynamic audio cues.
    \item \textbf{Save System:} Serializes structural parameters into JSON configurations.
\end{itemize}

\noindent Vertical movement computations utilize a constant gravity scalar to determine the instantaneous speed vector $v_y$ at a given time step $t$:
\begin{equation}
v_y = v_0 - (g \cdot t)
\end{equation}


\section{User Interface}

\subsection{How to Use (Controls)}
The input registration system translates physical inputs into unified logical controls. Table~\ref{tab:controls} highlights the primary hardware map allocations.

\begin{table}[h]
\centering
\caption{Hardware Controller Input Map Settings}
\label{tab:controls}
\begin{tabular}{|l|l|l|}
\hline
\textbf{Action Trigger} & \textbf{Primary Input (Keyboard)} & \textbf{Secondary Input (Gamepad)} \\ \hline
Move Left / Right       & A / D Keys                        & Left Analog Stick                  \\ \hline
Jump                    & Spacebar                          & Button South (A/Cross)             \\ \hline
Dash Impulse            & Left Shift                        & Right Trigger (R2/RT)              \\ \hline
\end{tabular}
\end{table}

\subsection{Database Explanation \& Data Management}
Persistent metrics, profile metrics, and leaderboard historical frames map into a localized SQLite instance. Level layout configuration properties operate entirely on text-based layouts created inside external level design suites.


% ==========================================
% v. GLOSSARY
% ==========================================
\section{Glossary}
\begin{table}[h]
\centering
\caption{Acronym Reference Guide}
\label{tab:glossary}
\begin{tabular}{|l|l|}
\hline
\textbf{Acronym} & \textbf{Definition} \\ \hline
\textbf{FPS}     & Frames Per Second \\ \hline
\textbf{FSM}     & Finite State Machine \\ \hline
\textbf{JSON}    & JavaScript Object Notation \\ \hline
\textbf{QA}      & Quality Assurance \\ \hline
\textbf{SDD}     & Software Design Document \\ \hline
\end{tabular}
\end{table}


% ==========================================
% vi. REFERENCES
% ==========================================
% Using explicit bibitems to align cleanly with standard Springer style
\begin{thebibliography}{8}

\bibitem{sasyr}
SASYR 2025 Submission Framework Guidelines, \url{http://sasyr.ipb.pt/}. Last accessed May 2026.

\bibitem{patterns}
Nystrom, R.: Game Programming Patterns: State Architectures. Genever Benning (2014). \url{https://gameprogrammingpatterns.com/state.html}.

\bibitem{lncs}
Springer Lecture Notes in Computer Science (LNCS) Authors Guidelines (2017).

\end{thebibliography}

\end{document}